\documentclass[11pt,a4paper]{article}
\usepackage[margin=2.3cm]{geometry}
\usepackage{amsmath,amssymb}
\usepackage{booktabs,longtable,array}
\usepackage[version=4]{mhchem}
\usepackage{xcolor}
\usepackage[colorlinks=true,linkcolor=blue!60!black,urlcolor=blue!60!black]{hyperref}
\usepackage{enumitem}
\usepackage{fancyhdr}
\pagestyle{fancy}\fancyhf{}
\rhead{GreenhouseTwin n79}\lhead{Model description}\cfoot{\thepage}
\newcommand{\code}[1]{\texttt{#1}}
\sloppy

\title{\textbf{GreenhouseTwin n79}\\[4pt]\large The greenhouse--cucumber model without flags:\\
climate, photosynthesis, crop growth and LAI}
\author{Greenhouses project --- model reference}
\date{Branch \code{n79}, 21 October 2026}

\begin{document}
\maketitle
\begin{abstract}
\noindent This document specifies the model coded in \code{invernaderosN79/GreenhouseTwin}: one equation set, no
behaviour flags. It is the configuration chosen as the baseline for making the physics more descriptive: the V2.2
engine with the n79 settings (fixed compartment geometry, energy-conserving solar column, cover node with
condensation, open-loop logged vents, stack leakage, closed floor balance, internal mass, full lamp heat, crop
respiration) and Aaron V\'elez's generic cucumber model with its canopy and LAI extension (the n105 crop set-up).
Every equation is written with the term names of the code (\code{src/climate.jl}, \code{src/crop.jl}), so each line can
be checked against \code{model\_terms}. Parameter tables are generated from the parameter files. Known structural
gaps are listed at the end; they are the starting point of the physics work.
\end{abstract}
\tableofcontents

%======================================================================
\section{Overview}\label{sec:overview}
%======================================================================
\paragraph{Scope.} One compartment of the Autonomous Greenhouse Challenge 2018 (AiCU, Bleiswijk; 96\,m$^2$, Venlo, gutter
5\,m), cucumber crop transplanted on 14 August 2018, 115 simulated days. All quantities are per m$^2$ of floor.

\paragraph{Two models, offline coupling.}
\begin{itemize}[nosep]
\item The \textbf{climate model} is a system of 8 ODEs integrated day by day (Tsit5, 5-min output):
\begin{equation}
u=\bigl[\,T_1,\;T_2,\;C_1,\;V_1,\;e_{\rm int},\;T_5,\;T_d,\;T_m\,\bigr],
\end{equation}
canopy temperature, air temperature, \ce{CO2} concentration [mg\,m$^{-3}$], vapour pressure [Pa], integral of the heating error,
floor layer, deep soil and internal mass temperatures [K].
\item The \textbf{crop model} is a daily step (Section~\ref{sec:crop}) driven by measured climate. Its daily LAI, non-leaf
respiration and standing fresh mass are inputs of the climate model (\code{test/crop\_series.jl}). The crop does not see
the simulated climate (live coupling is a planned extension).
\end{itemize}

\paragraph{Provenance and verification.} The code is GreenhouseSim V2.2 (git tag \code{v2.2-flags-final}) with every flag
resolved to its n79 value and dead branches removed. \code{test/equivalence\_synthetic.jl} and
\code{test/equivalence\_n79.jl} compare it with V2.2 run with the n79 flags (Section~\ref{sec:verify}).

\paragraph{Inputs.} Table~\ref{tab:inputs}. Weather and LAI are functions of absolute time; controls and crop inputs of
time since the season start. All are 5-min step functions.
\begin{table}[h]\centering\small
\caption{Model inputs (\code{src/inputs.jl}; AGC-2018 construction in \code{src/agc\_data.jl}).}\label{tab:inputs}
\begin{tabular}{>{\raggedright\arraybackslash}p{0.17\textwidth}>{\raggedright\arraybackslash}p{0.24\textwidth}>{\raggedright\arraybackslash}p{0.5\textwidth}}
\toprule symbol & code & meaning (AiCU source)\\\midrule
$I_{\rm glob}$ & \code{Iglob} & global radiation outside [W\,m$^{-2}$] (meteo)\\
$T_{\rm sky}$ & \code{Tsky} & sky temperature from the pyrgeometer [K]\\
$T_o$, $v$, $V_o$ & \code{Tout, Wind, VPout} & outside temperature, wind speed, vapour pressure\\
$C_o$, $T_{\rm mc}$ & \code{CO2out, TmechCool} & outside \ce{CO2} 834.7\,mg\,m$^{-3}$; mechanical-cooling surface 273.15\,K\\
$L$ & \code{LAI} & leaf area index per m$^2$ floor (crop series)\\
$U_1$, $U_{\rm bl}$ & \code{EnScr, BlackScr} & energy and blackout screen closure\\
$U_{\rm lee}$, $U_{\rm win}$ & \code{VentLee, VentWind} & leeward, windward roof-vent openings\\
$U_{10}$ & \code{CO2Valve} & \ce{CO2} valve opening (setpoint rule scaled to the logged daily dosage)\\
$U_{11}$ & \code{HeatSP} & heating setpoint [K]\\
$U_{12}$ & \code{Lamps} & lamp fraction\\
$T_{\rm log}$ & \code{Tpipe} & logged pipe temperature [K] (both circuits folded; no reading $\to T_o+2$)\\
$R_{\rm crop}$, FW & \code{CropResp, CropFW} & crop respiration at 20\,$^\circ$C [mg\,\ce{CO2}\,m$^{-2}$\,s$^{-1}$]; standing fresh mass [kg\,m$^{-2}$]\\
$U_{2\ldots7}$, $U_9$ & \code{FanPad}, \code{MechCool}, \code{Heater}, \code{Shade}, \code{SideVent}, \code{ForcedVent}, \code{Fog} & other actuators; zero in the AiCU season\\
\bottomrule\end{tabular}
\end{table}

%======================================================================
\section{Climate model}\label{sec:climate}
%======================================================================
\subsection{State equations}
With $c_{\rm air}=\phi_2\rho_3\alpha_5$ the air heat capacity per m$^2$ floor and $m_{\rm air}$ = \code{air\_cap\_mult},
\begin{align}
\alpha_1 L\,\dot T_1 &= r_1 + Q_{lc} + r_6 - h_1 - l_1 - r_7 - r_{14},\label{eq:T1}\\
m_{\rm air}\,c_{\rm air}\,\dot T_2 &= h_1 + h_2 + h_4 - h_7 - h_{11} - r_{12} - r_{13} - l_2 + Q_{sa} - Q_{af} - Q_{am} + Q_{la},\label{eq:T2}\\
\phi_2\,\dot C_1 &= o_1 + o_2 + o_3 - A - o_5 + o_R,\label{eq:C1}\\
\kappa_3\,\dot V_1 &= p_1 + p_2 + p_3 + p_4 - p_5 - p_6 - p_7 - p_8 - p_9 - p_{10},\qquad \kappa_3=\frac{\psi_1\phi_2}{\omega_2T_2},\label{eq:V1}\\
\dot e_{\rm int} &= \text{(anti-windup integral of } U_{11}-T_2\text{, eq.~\ref{eq:pi})},\\
C_{\rm soil}\,\dot T_5 &= Q_{af} - Q_{fd} + h_{11} + r_{14} - r_{11} + Q_{fs},\label{eq:T5}\\
\dot T_d &= (T_5-T_d)/\tau_d,\\
\bigl(C_m + c_{p,\rm crop}\,\mathrm{FW}\bigr)\,\dot T_m &= Q_{am} + Q_{ms}.\label{eq:Tm}
\end{align}
The cover temperature $T_c$ and the screen temperature $T_{\rm scr}=\tfrac12(T_o+T_2)$ are algebraic. Every term is defined
below with the name it has in \code{model\_terms}.

\subsection{Solar radiation: the energy-conserving column}\label{sec:solar}
For each band $b$ (PAR $I_P=\eta_2I_{\rm glob}$, NIR $I_N=\eta_3I_{\rm glob}$, $\eta_2+\eta_3=1$) five elements are stacked
from the sky down: roof, energy screen, blackout screen, canopy, opaque floor. Each layer is symmetric with transmission
$t$ and reflection $r$:
\begin{align}
\text{roof:}\quad & t_r=(1-\eta_1)\tau_1\ (=0.65),\qquad r_r=1-t_r-\eta_1-a_c\ (=0.15),\\
\text{screen at closure }U:\quad & t=1-U(1-\tau_s),\qquad r=U\rho_s\qquad\text{(Vanthoor 2011, 8.14--8.17)},\\
\text{canopy:}\quad & t_c=e^{-kL},\qquad r_c=\rho_\infty(1-t_c)\qquad\text{(Vanthoor 8.31--8.32)},
\end{align}
with $k=\beta_1$, $\rho_\infty=\rho_1$ (PAR) and $k=\beta_3$, $\rho_\infty=$ \code{rho\_can\_nir} (NIR); floor reflectance
$\rho_2$ (PAR), $1-\alpha_7$ (NIR); screens: energy $(\tau_s,\rho_s)=(0.75,0.15)$, blackout $(0.02,0.50)$, NIR = PAR
(provisional). With all repeated reflections (adding method), from the bottom up,
\begin{align}
R_4&=r_4+\frac{t_4^2\rho_f}{1-r_4\rho_f},\quad R_j=r_j+\frac{t_j^2R_{j+1}}{1-r_jR_{j+1}}\ (j=3,2,1),\\
D_2&=\frac{t_1}{1-r_1R_2},\quad D_{j+1}=\frac{t_jD_j}{1-r_jR_{j+1}}\ (j=2,3),\quad D_5=\frac{t_4D_4}{1-r_4\rho_f},\\
A_1&=(1-t_1-r_1)(1+R_2D_2),\quad A_j=(1-t_j-r_j)(D_j+R_{j+1}D_{j+1}),\quad A_4=(1-t_4-r_4)(D_4+\rho_fD_5),\\
A_f&=(1-\rho_f)D_5,\qquad R_1+A_1+A_2+A_3+A_4+A_f=1 .
\end{align}
Booking [W\,m$^{-2}$]: canopy $r_1=Q_{\rm can}=I_PA_4^P+I_NA_4^N$; floor level $Q_{\rm flr}=I_PA_f^P+I_NA_f^N$, of which the
floor gets $Q_{fs}=(1-f_m)Q_{\rm flr}$ and the internal mass $Q_{ms}=f_mQ_{\rm flr}$; roof absorption
$Q_{\rm roof}=\sum_bI_bA_1^b$ splits into structure $Q_{st}=\sum_bI_bA_1^b\min(\eta_1/(1-t_r-r_r),1)$ (to the air) and glass (to the
cover); screen absorption $Q_{\rm scr}=\sum_bI_b(A_2^b+A_3^b)$ goes \code{scr\_abs\_air} to the air and the rest to the cover:
\begin{equation}
Q_{sa}=Q_{st}+s_aQ_{\rm scr},\qquad Q_{\rm cov}=(Q_{\rm roof}-Q_{st})+(1-s_a)Q_{\rm scr},\qquad
Q_{\rm can}+Q_{\rm flr}+Q_{sa}+Q_{\rm cov}+Q_{\rm out}=I_{\rm glob}.
\end{equation}

\subsection{Plant light and photosynthesis}\label{sec:photo}
Stomata and photosynthesis see the PAR on both canopy faces plus lamp PAR:
\begin{equation}
\mathrm{PAR}_p = I_PD_4^P + I_P\rho_2D_5^P + P_{\rm lamp}\,\mathrm{clamp}(U_{12},0,1)\qquad[\text{W\,m}^{-2}],
\end{equation}
($P_{\rm lamp}$ = \code{lamp\_par\_W}). Canopy assimilation (FvCB, constant effective conductance $g=g_sg_m/(g_s+g_m)$) is
\begin{equation}
A = 44010\;\frac{1-e^{-k_pL}}{k_p}\;\min(A_j,A_c,A_p)\quad[\text{mg\,\ce{CO2}\,m}^{-2}\text{s}^{-1}],\qquad k_p=0.7,
\end{equation}
evaluated at $T_2$, $C_1$, $\mathrm{PAR}_p$ with:
\begin{align}
I&=4.6\times10^{-6}\,\mathrm{PAR}_p,\quad C_i=0.509\times10^{-6}C_1,\quad
\Gamma^*=\frac{0.5\,O_a}{S_{c/o,25}\exp\bigl(\frac{T-298.15}{T}\frac{E_{Soc}}{298.15R}\bigr)},\\
J&=\frac{(\alpha I+J_{\max})-\sqrt{(\alpha I+J_{\max})^2-4\theta\alpha IJ_{\max}}}{2\theta},\quad
V_{c\max}=\frac{V_{c\max,25}\,Q_{10,V}^{0.1(T-298.15)}}{1+e^{0.128(T-315.15)}},\\
K_C&=K_{C25}Q_{10,KC}^{0.1(T-298.15)},\quad K_O=K_{O25}Q_{10,KO}^{0.1(T-298.15)},\quad R_d=R_{d,25}\,2^{(T_1-298.15)/10},
\end{align}
and $A_c$, $A_j$ the smaller roots of
\begin{align}
A_c^2-\bigl(V_{c\max}+g(C_i+K_C(1+O_a/K_O))-R_d\bigr)A_c+g\bigl(V_{c\max}(C_i-\Gamma^*)^+-(C_i+K_C(1+O_a/K_O))R_d\bigr)&=0,\\
A_j^2-\bigl(\tfrac14J-R_d+g(C_i+2\Gamma^*)\bigr)A_j+\tfrac14gJ(C_i-\Gamma^*)^+-g(C_i+2\Gamma^*)R_d&=0,
\end{align}
($-R_d$ if the discriminant is negative) and $A_p=1.5V_{c\max}-R_d$.

\subsection{Stomata and transpiration}
With $q_2=P_s(T_1)$ the saturated vapour pressure at the canopy and the day/night switch
$q_7=\bigl(1+e^{(\mathrm{PAR}_p-\delta_1)/\gamma_5}\bigr)^{-1}$ ($\approx1$ by day),
\begin{align}
q_8&=\delta_5+(\delta_4-\delta_5)q_7,\quad q_9=\delta_7+(\delta_6-\delta_7)q_7,\quad
q_{10}=\frac{\mathrm{PAR}_p+\delta_2}{\mathrm{PAR}_p+\delta_3},\\
q_3&=\gamma_4\,q_{10}\,\bigl[1+q_8(C_1T_2/536.4-200)^2\bigr]\bigl[1+q_9(q_2-V_1)^2\bigr]\quad\text{(stomatal resistance, CO}_2\text{ in ppm)},\\
q_1&=\frac{2\rho_3\alpha_5L}{\gamma\gamma_2(\gamma_3+q_3)}\quad(q_2\ge V_1),\qquad
q_1=\frac{2\rho_3\alpha_5L}{\gamma\gamma_2\gamma_3}\quad(q_2<V_1,\ \text{dew}),
\end{align}
transpiration $p_1=q_1(q_2-V_1)$ [kg\,m$^{-2}$\,s$^{-1}$] and its latent heat $l_1=\gamma_2p_1$.

\subsection{Cover node}\label{sec:cover}
The cover has no heat capacity. Linearising its long-wave exchange with the crop air and the sky,
$k_R=4\varepsilon_3\sigma\bigl(\tfrac12(T_2+T_{\rm sky})\bigr)^3$,
\begin{equation}
T_c^{(0)}=\frac{k_R(T_2+T_{\rm sky})+h_{\rm out}T_o+Q_{\rm cov}}{2k_R+h_{\rm out}},\qquad
T_c=T_c^{(0)}+\frac{\gamma_2\,f_{\rm cov}\,6.4\!\times\!10^{-9}\cdot1.7\,|T_2-T_c|^{0.33}\,(V_1-P_s(T_c))^+}{2k_R+h_{\rm out}},
\end{equation}
the second equation solved by four fixed-point passes (latent heat of the vapour condensing on the cover warms it), with
$f_{\rm cov}=1-U_1(1-p_{\rm scr})$ the share of crop air in contact with the cover ($p_{\rm scr}$ = \code{vent\_scr\_perm}).

\subsection{Long-wave radiation}\label{sec:lw}
With $g_1=0.49(1-e^{-\beta_3L})$, $g_2=1-U_1(1-\tau_3)$, $g_3=\tau_1g_2(1-0.49\pi\gamma_1\phi_1)e^{-\beta_3L}$,
$F=\varepsilon_3\sigma(T_2^4-T_c^4)$:
\begin{align}
r_6&=\alpha_3\varepsilon_1\varepsilon_2\,g_1\,\sigma(T_p^4-T_1^4) &&\text{pipe}\to\text{canopy}\\
r_7&=\frac{g_1}{0.49}\varepsilon_2\varepsilon_3\,g_2\,\sigma(T_1^4-T_c^4) &&\text{canopy}\to\text{cover}\\
r_{11}&=\varepsilon_4\,g_3\,\varepsilon_3\sigma(T_5^4-T_c^4) &&\text{floor}\to\text{cover}\\
r_{12}&=\varepsilon_{\rm scr}\,U_1\tau_2\,F &&\text{energy screen (at air }T)\to\text{cover}\\
r_{13}&=A_{\rm cov}\,\varepsilon_6\,F &&\text{air}\to\text{cover}\\
r_{14}&=\varepsilon_4\varepsilon_2(1-0.49\pi\gamma_1\phi_1)\,\sigma(T_1^4-T_5^4) &&\text{canopy}\to\text{floor}
\end{align}

\subsection{Air exchange}\label{sec:vent}
Roof vents: the leeward and windward sides are combined into one equivalent opening,
\begin{align}
B&=\omega_1\nu_6\frac{T_2-T_o}{T_2+T_o},\quad w=n_3v^2,\quad n_3=\nu_2(1-\eta_{11}U_5),\quad
\Phi=\tfrac12\Bigl[U_{\rm lee}^{\,e}\sqrt{(B+c_lw)^+}+g_wU_{\rm win}^{\,e}\sqrt{(B+c_ww)^+}\Bigr],\\
U_8&=\Phi/\sqrt{(B+w)^+}\quad\bigl(\text{if }(B+w)^+>10^{-12};\ \text{else }\tfrac12(U_{\rm lee}^e+g_wU_{\rm win}^e)\bigr),
\end{align}
($e,g_w,c_l,c_w$ = \code{vs\_exp, vs\_gwind, vs\_clee, vs\_cwind}). Air exchange rates [m$^3$\,m$^{-2}$\,s$^{-1}$]:
\begin{align}
f_7&=\frac{U_8\,f_{\rm scr}\,\nu_5n_1}{2\alpha_6}\sqrt{(B+w)^+},\quad f_{\rm scr}=1-U_1(1-p_{\rm scr}),\quad n_1=\nu_1(1-\eta_{10}U_5) &&\text{roof vents}\\
f_4&=\eta_6f_7,\qquad f_2=\eta_6U_1f_5,\quad f_5=\frac{n_1\nu_3U_6\,v\sqrt{n_3}}{2\alpha_6} &&\text{roof, side vents}\\
f_3&=U_7\phi_8/\alpha_6 &&\text{forced ventilation}\\
f_6&=\nu_4\sqrt{\max(v,0.25)^2+c_s|T_2-T_o|} &&\text{leakage (wind + stack)}
\end{align}
The same total $f_2+f_3+f_4+f_6$ carries heat, vapour and \ce{CO2}.

\subsection{Heating}\label{sec:heat}
Where the logged pipe is warm ($T_{\rm log}>T_2+1$) it is used as is; elsewhere a PI controller requests
\begin{equation}\label{eq:pi}
e=U_{11}-T_2,\quad \Delta=K_pe+K_i\max(e_{\rm int},0),\quad T_p=\max\bigl(T_{\rm log},\,T_2+\mathrm{clamp}(\Delta,0,T_{\max}-T_2)\bigr),
\end{equation}
and $\dot e_{\rm int}=e$ except when the logged pipe is warm, the request saturates with $e>0$, or $e_{\rm int}\le0$ with $e<0$ (then 0).
Pipe-to-air convection: $h_4=n_p\,1.99\pi\phi_1\gamma_1|T_p-T_2|^{0.32}(T_p-T_2)$.

\subsection{Lamps}
The full electrical power $Q_l=P_{\rm el}U_{12}$ is released inside: $Q_{lc}=f_{\rm can}(1-e^{-\beta_1L})Q_l$ to the canopy,
$Q_{la}=Q_l-Q_{lc}$ to the air ($P_{\rm el}=117$\,W\,m$^{-2}$, $f_{\rm can}$ = \code{lamp\_fcan}).

\subsection{Sensible and latent heat of the air}
\begin{align}
h_1&=2\alpha_4L(T_1-T_2) &&\text{canopy}\to\text{air}\\
h_2&=\frac{U_2\phi_7}{\alpha_6}\bigl(\rho_3\alpha_5T_o-\gamma_2\rho_3\eta_5(\phi_5-\phi_6)\bigr) &&\text{fan-pad}\\
h_7&=\rho_3\alpha_5(f_2+f_4+f_3+f_6)(T_2-T_o) &&\text{ventilation and leakage}\\
h_{11}&=2\nu_7(T_2-T_5)/(\phi_2+\nu_8) &&\text{air}\to\text{floor conduction}\\
Q_{af}&=k_{\rm ground}(T_2-T_5),\quad Q_{am}=h_m(T_2-T_m) &&\text{air}\to\text{floor, mass convection}\\
l_2&=\gamma_2U_9\phi_9/\alpha_6 &&\text{fog evaporation}
\end{align}

\subsection{Vapour}
\begin{align}
p_2&=\rho_3\frac{U_2\phi_7}{\alpha_6}\bigl(\eta_5(\phi_5-\phi_6)+\phi_6\bigr),\quad p_3=\frac{U_9\phi_9}{\alpha_6},\quad
p_4=\frac{\eta_{12}U_4\lambda_4}{\alpha_6},\quad p_6=\frac{U_2\phi_7}{\alpha_6}\frac{\psi_1}{\omega_2}\frac{V_1}{T_2},\\
p_5&=\frac{\psi_1}{\omega_2}\Bigl(\frac{V_1}{T_2}-\frac{V_o}{T_o}\Bigr)(f_2+f_3+f_4+f_6),\qquad
p_7=U_3\,6.4\!\times\!10^{-9}(V_1-P_s(T_{\rm mc}))^+,\\
p_8&=6.4\!\times\!10^{-9}\cdot1.7\,U_1|T_2-T_{\rm scr}|^{0.33}(V_1-P_s(T_{\rm scr}))^+ &&\text{screen condensation}\\
p_9&=6.4\!\times\!10^{-9}\cdot1200\,|T_1-T_2|^{0.33}(V_1-P_s(T_2))^+ &&\text{supersaturated air}\\
p_{10}&=A_{\rm cov}f_{\rm cov}\,6.4\!\times\!10^{-9}\cdot1.7\,|T_2-T_c|^{0.33}(V_1-P_s(T_c))^+ &&\text{cover condensation}
\end{align}

\subsection{\ce{CO2}}
\begin{equation}
o_1=\frac{\eta_{13}U_4\lambda_4}{\alpha_6},\ o_2=\frac{U_{10}\psi_2}{\alpha_6},\ o_3=\frac{U_2\phi_7}{\alpha_6}(C_1-C_o),\
o_5=(C_1-C_o)(f_2+f_3+f_4+f_6),\ o_R=R_{\rm crop}\,2^{(T_1-293.15)/10}.
\end{equation}

\subsection{Floor, deep soil and internal mass}
$Q_{fd}=k_{\rm deep}(T_5-T_d)$; the floor balance (\ref{eq:T5}) collects air convection and conduction, canopy long-wave,
its own long-wave loss to the cover and its share of the solar column. The internal mass (structure, slabs, water, standing
crop) exchanges $Q_{am}$ with the air and absorbs $Q_{ms}$.

\subsection{Psychrometrics}
$P_s(T)=611.21\exp\bigl[(18.678-\frac{T_c}{234.5})\frac{T_c}{257.14+T_c}\bigr]$ above 0\,$^\circ$C (ice form below),
$T_c=T-273.15$; $\mathrm{RH}=100V_1/P_s(T_2)$.

%======================================================================
\section{Crop model: Aaron V\'elez's generic cucumber model with canopy and LAI}\label{sec:crop}
%======================================================================
CROPSECTION

%======================================================================
\section{Parameters}\label{sec:params}
%======================================================================
PARAMTABLES

%======================================================================
\section{Parameter inference}\label{sec:inference}
%======================================================================
INFERENCE

%======================================================================
\section{Known structural gaps}\label{sec:gaps}
%======================================================================
GAPS

%======================================================================
\section{Verification}\label{sec:verify}
%======================================================================
VERIFY

\end{document}
